\documentclass[xcolor=svgnames]{beamer}
\usepackage{latexsym,amssymb,amsmath,stmaryrd}
\usepackage[utf8]{inputenc}
\usepackage[T1]{fontenc}
\usepackage[francais]{babel}
\usepackage{pgf}
\usepackage{tikz}
\usetikzlibrary{arrows,patterns,plotmarks,shapes,snakes,er,3d,automata,%
backgrounds,topaths,trees,petri,mindmap}
\usepackage{pgfbaseimage}
\usepackage{ifpdf}
\ifpdf
\usepackage{graphicx}
\else
%\usepackage[dvips]{graphicx}
\usepackage{pstricks,pst-tree,pst-node}
\newpsobject{showgrid}{psgrid}{subgriddiv=1,griddots=10,gridlabels=6pt}
\fi
\usepackage{moreverb}
\usepackage[normalem]{ulem}
\usepackage{version}
\usepackage{url}
\usepackage{multicol}
\usepackage{listings}
\lstset{
  language=[ANSI]C, 
%  gobble=2, 
%  escapeinside="", 
  basicstyle=\ttfamily, 
%  directivestyle=\color{Fuchsia},
  identifierstyle = \color{DarkOrange},
  keywordstyle=\color{DarkGreen}, 
  commentstyle=\color{red}, 
%  numbers=left, 
%  numbersep=5pt, 
%  numberstyle=\scriptsize, 
%  moredelim=[is][\color{gray}\itshape]{/*}{*/}, 
%  moredelim=[is][\alert]{/+}{+/}, 
%  morecomment=[is]{/=}{=/}, 
%  morecomment=[in]{comment=}{,} 
%  fancyvrb=true 
} 

\newcommand{\binaire}[1]{\ensuremath{\underline{#1}}}
\newcommand{\C}[1]{\texttt{#1}}
\newcommand{\bbbn}{\ensuremath{\mathbb{N}}}
\input{communs/cc_beamer}
\usetheme{classic}
\newcommand{\nowrite}{\put(10,-4){\includegraphics[scale=.05]{creative_commons/nopencil}}}
\newcommand{\youwrite}{\put(10,-4){\includegraphics[scale=.05]{creative_commons/pencil}}}
\newcommand{\aemporter}{\put(10,-6){\includegraphics[scale=.05]{creative_commons/szymonraj_Shopping_bag}}}


%%% Titre -- Cours 1
\title{Éléments d'informatique -- Cours 1. Éléments d'architecture des
ordinateurs, mini-assembleur}
\author{Pierre Boudes}
\date{\today}

\begin{document}

%% Page de titre et licence CC.
\begin{frame}
	\titlepage
	\vfill
	\begin{center}
		\CcGroupByNcSa{0.83}{0.95ex}\\[2.5ex]
		{\tiny\CcNote{\CcLongnameByNcSa}}
		\vspace*{-2.5ex}
	\end{center}
\end{frame}

\section[Plan]{}
\frame[label=plan]{\tableofcontents}

  \section[Plan du cours]{Survol du contenu du cours (ce semestre)}
\begin{frame}
  \frametitle{Survol du contenu du cours (ce semestre)\nowrite}
\pause
  \begin{itemize}
  \item Éléments d'architecture des ordinateurs (+mini-assembleur)\pause
  \item Éléments de systèmes d'exploitation\pause
\item Programmation structurée impérative (éléments de langage C)\pause
\begin{itemize}
  \item Structure d'un programme C
  \item Variables : déclaration (et initialisation), affectaction
  \item Évaluation d'expressions\pause
  \item Instructions de contrôle : if, for, while\pause
  \item Types de données : entiers, caractères,
    réels, tableaux, types composés (enregistrements)\pause
  \item Fonctions d'entrées/sorties (scanf/printf)\pause
  \item Écriture et appel de fonctions\pause
  \item Débogage\pause
\end{itemize}
\item Notions de compilation
  \begin{itemize}
  \item Analyse lexicale, analyse syntaxique, analyse sémantique
  \item préprocesseur du compilateur C (include, define)
   \item Édition de lien\pause
  \end{itemize}
\item Algorithmes élémentaires\pause
\item Méthodologie de résolution, manipulation sous linux
\end{itemize}
\end{frame}


\begin{frame}
  \begin{quote}
  « L'informatique n'est pas plus la science des ordinateurs que
  l'astronomie n'est celle des télescopes. »\\ 
  \hfill E.~W.~Dijkstra\\[1cm]
\end{quote}
\end{frame}

\section{Architecture de von Neumann}
\begin{frame}
  \frametitle{Architecture de von Neumann}
  \begin{itemize}
  \item John William Mauchly et John Eckert autant (ou plus) que
    vN\pause
  \item Qu'est-ce que c'est ?\pause
    \begin{itemize}      
  \item L'idée d'une machine à \alert{programme
      stocké}\pause
\item  Une machine réalisée, l'ancêtre de nos processeurs\pause
\end{itemize}
  \item De quoi  cette machine est-elle faite ?\pause
    \begin{itemize}
    \item De mémoire (une suite de cases numérotées)\pause
    \item d'une \alert{unité de calcul}, travaillant sur des \alert{registres}\pause
    \item d'un \alert{bus} système (adresses et données) reliant mémoire
      et UC\pause
    \item \sout{De périphériques} (on oublie !)\pause
    \item La mémoire contient le programme et les données.
    \end{itemize}
  \end{itemize}
\end{frame}


\section{Représentation  des informations}

\begin{frame}
  \frametitle{Représentation en binaire des informations}
  \pause
  \begin{definition}[bit]
    \begin{itemize}
    \item Le chiffre binaire, ou \emph{bit}, est l'équivalent binaire
      de nos chiffres décimaux. Il peut valoir soit $0$ soit $1$. Un
      bit est une \alert{quantité élémentaire d'information} (oui ou non,
      ouvert ou fermé, etc.).\pause
    \item L'information manipulée par un ordinateur est consituée de
      bits.\pause
\item Les cases mémoires et les registres contiennent des \alert{mots mémoire}
  : des suite de $n$ bits, où $n$ est fixé une fois pour toute par
  l'architecture matérielle.\pause
\item les instructions du langage machine sont écrites en binaire.\pause
\item   le \alert{langage assembleur} est une notation du langage
  machine plus pratique pour les humains. 
   \end{itemize}
  \end{definition}
\pause Nous en verrons un peu plus sur les codages en binaire des données
dans un autre cours.
\end{frame}

\section{Cycle d'exécution}
\begin{frame}
  \frametitle{Cycle d'exécution}
\pause
\begin{itemize}
  \item Le registre \alert{compteur de programme} (CP) contient l'adresse du mot mémoire
    représentant la prochaine instruction\pause
\item le contenu de ce mot est transféré de la mémoire centrale dans
  le \alert{registre d'instruction} (RI)\pause
\item CP est \emph{incrémenté} (c'est à dire que sa valeur augmente de 1)\pause
\item le contenu de RI est décodé afin de déterminer l'opération à
  exécuter\pause
\item l'opération est exécutée (le contenu d'un ou plusieurs
  registres est modifié, ou bien celui d'une case mémoire)\pause
\item Fin du cycle d'exécution et démarrage d'un nouveau cycle
\end{itemize}
\end{frame}

\section{Instructions}
\begin{frame}
  \frametitle{Instructions}
Une instruction type comporte un \alert{code d'opération} et, si nécessaire, une
ou deux \emph{opérandes} (ou \emph{arguments} de l'opération).

\pause
\begin{block}{Vocabulaire} Dans l'expression arithmétique usuelle $3 + 5$,  le signe $+$ est l'opérateur
et les nombres $3$ et $5$ sont les opérandes.
\end{block} 
\end{frame}

\begin{frame}
\frametitle{Quelques instructions typiques (Amil)\nowrite}
\pause
  \begin{tabular}[c]{lp{8.3cm}}
  \C{stop} & Arrête l'exécution du programme.\\ 
  \C{noop} & N'effectue aucune opération.\\ \pause
  \C{lecture i rj} & Charge, dans le registre $j$, le contenu de la mémoire d'adresse $i$.\\
  \C{ecriture ri j} & Écrit le contenu du registre $i$ dans la mémoire
  d'adresse $j$.\\ \pause
  \C{saut i} & Met CP à la valeur $i$.\\
  \C{sautpos ri j} & Si la valeur contenue dans le registre $i$ est
  positive ou nulle, met CP à la valeur $j$.\\ \pause
%  \C{valeur x ri} & Initialise le registre $i$ avec la valeur $x$.\\
 \C{inverse ri} & Inverse le signe du contenu du registre $i$.\\ 
  \C{add ri rj} & Ajoute la valeur du registre $i$ à celle du registre $j$.\\
  \C{soustr ri rj} & Soustrait la valeur du registre $i$ à celle du registre $j$.\\
  \C{mult ri rj} & Multiplie \ldots\\
  \C{div ri rj} & Divise \ldots \\ \pause
% \\
%   \multicolumn{2}{l}{\textbf{Instructions plus avancées}}\\
%   \C{et ri rj} & Effectue le et bit à bit de la valeur du registre
%   $i$ et de celle du registre $j$. Le résultat est placé dans le
%   registre $j$.\\
   \C{lecture *ri rj} & Charge, dans $rj$, le contenu de la mémoire dont l'adresse est la valeur du registre $i$\\
%   \C{ecriture ri *rj} & Écrit le contenu de $ri$ dans la  mémoire dont  l'adresse est la valeur du registre $j$.
\end{tabular}
\end{frame}


\section{Trace d'exécution}
\begin{frame}[fragile]
\frametitle{Trace d'exécution\nowrite}
On simule pas à pas l'exécution du programme.\pause
\begin{columns}[t]
  \begin{column}{3.5cm}
\newcommand{\N}[1]{\structure{\scriptsize #1.}}
\pause \structure{Programme}
\begin{semiverbatim}
\N{ 1} lecture 10 r0
\N{ 2} lecture 11 r2
\N{ 3} soustr r2 r0
\N{ 4} sautpos r0 8
\N{ 5} lecture 10 r2
\N{ 6} add r2 r0
\N{ 7} saut 4
\N{ 8} ecriture r0 12
\N{ 9} stop
\N{ 10} 14
\N{ 11} 5
\N{ 12} ?
\end{semiverbatim}
 \end{column}
  \begin{column}{10cm}
    \pause \structure{Trace}\\[0.5cm]
\scriptsize
\newcommand{\commentaire}[1]{\pause} %hack
\begin{tabular}[c]{l||c|c|c|c|c|c|c|}
\hline
 \emph{Instructions} & Cycles & CP& r0& r2& 10& 11 &12\\ \hline \pause
\hfill Initialisation & 0 & 1 & ? & ? & 14 & 5 & ?
 \\ \hline \commentaire{Lecture de la donnée d'adresse 10 dans le
   registre 0} 
\C{lecture 10 r0} & 1 & 2  & 14 & & & &\\ \hline
 \commentaire{Initialisation du registre 2 à 5
} \C{lecture 11 r2} & 2 & 3  & & 5 & & &\\ \hline
 \commentaire{Soustrait la valeur du registre 2 au registre 0
} \C{soustr r2 r0
} & 3 & 4  & 9 & & & &\\ \hline
 \commentaire{Si la valeur (9) du registre 0 est positive, saute à l'adresse 8
} \C{sautpos r0 8
} & 4 & \textbf{8} & & & & &\\ \hline
 \commentaire{Écriture du registre 0 à l'adresse 12
} \C{ecriture r0 12
} & 5 & 9  & & & & & 9
\\ \hline
 \commentaire{Fin du processus.
} \C{stop
} & 6 & 10  & & & & &\\ \hline
\end{tabular}
  \end{column}
\end{columns}

\end{frame}

\section[Démos]{Démos et fin}
\begin{frame}
  \frametitle{Démos et fin}
\end{frame}
\end{document}
